% !TEX root = machine_learning.tex
\chapter{Boosting methods}

Boosting methods refer to algorithms in which classifiers are enhanced by
recursively training them on harder data. We first address the issue
of classification, and describe the original  algorithm (Adaboost). We then
interpret it as greedy gradient descent to optimize a combination of
classifiers.

\section{Classification}
Boosting for classification is essentially restricted to two-class
problems. In this context, we want to design a function $x \mapsto f(x) \in \{-1, 1\}$ on the
basis of a training set $(y_1, x_1), \ldots, (y_N, x_N)$. With the 0-1
loss, minimizing the  empirical error is equivalent to maximizing
$$
\CE_T(f) = \frac{1}{N} \sum_{i=1}^N y_i f(x_i)
$$

Boosting algorithms build a function $f$ as a weighted combination of
simple classifiers, $f_1, \ldots, f_M$, taking
$$
f(x) = \text{sign}\left(\sum_{k=1}^M \al_k f_k(x)\right)\,.
$$

Each $f_j$ is learned on the basis of weighted
data. We therefore assume that the  learning algorithm for simple classifiers takes as input the training set $X_1,\ldots,X_N$ as well as positive weights $w_1, \ldots, w_N$,   minimizing, for example, a loss function of the form
$$
\sum_{k=1}^N w_k R(y_k, f(x_k)).
$$
Boosting algorithms  estimate new $f_j$ and $\al_j$ at step $j$, while updating the current data weights. The basic algorithm works as follows.
\begin{enumerate}[label=(\arabic*),wide]
\item Start with uniform weights, letting  $w_k = 1/N$, $k=1, \ldots, N$.
\item Step $j$: fit a classifier $f_j$ using the current weights and
estimate the weighted error
$$\CE_{w} = \sum_{k=1}^N w_k \chi_{[y_k\neq f_j(x_k)]}/\sum_{k=1}^N
w_k.$$
Define $\al_j = (1/2)\log((1-\CE_w)/\CE_w)$ and update the weights by
$$ w_k \mapsto w_k \exp(-\al_j \chi_{[y_k\neq f_j(x_k)]}).$$
\item Iterate the previous step while $j\leq M$.
\end{enumerate}

Notice that for the algorithm to run, $\al_j$ should be positive (i.e., the $j$th classifier must do
better than pure chance on the training set).
The boosting effect comes from the fact that misclassified data at step $j$
are provided with a higher weight for the next step. This implies that
the algorithm focuses on the hardest cases. 


We now provide an interpretation of this algorithm. Even if the $f_q$ are binary, their linear combination is not
anymore. The classification rule only uses the sign of the product $Yf(X)$, but the value also has an
important interpretation, since, when it is positive, it can be interpreted as a margin with which
$X$ is correctly classified (provided, of course that $f$ is properly
normalized, since one can always improve the margin by multiplying $f$
by a constant). 

Assume that the function $f$ is evaluated, not only on the basis of
its classification error, but also based on this margin,
using a loss function of the kind
$$
\sum_{k=1}^N g(y_kf(x_k))
$$
where $g$ is decreasing. The boosting algorithm can then be interpreted as an classifier which
incrementally improves this cost function.
Denote by
$f^{(j)}$ the  classifier at step $j$:
$$
f^{(j)} = \sum_{q=1}^j \al_q f_q\,.
$$
The new classifier $f^{(j+1)}$ takes the form $f^{(j)} + \al_{j+1}
f_{j+1}$, where  
$f_{j+1}$ and $\al_{j+1}$ minimize, with respect
to $\tilde f$ and $\al \geq 0$, the function
$$
U(\al, \tilde f) = \sum_{k=1}^N g\left(y_kf^{(j)}(x_k) + \al y_k \tilde
f(x_k)\right)
$$

Using the fact that $\tilde f$ is a binary classifier, this can be written
\begin{eqnarray*}
U(\al, \tilde f) &=& \sum_{k=1}^N g(y_kf^{(j)}(x_k) + \al) \chi_{[y_k
= \tilde f(x_k)]} + \sum_{k=1}^N g(y_kf^{(j)}(x_k) - \al) \chi_{[y_k
\neq \tilde f(x_k)]}\\
&=& \sum_{k=1}^N (g(y_kf^{(j)}(x_k) - \al) - g(y_kf^{(j)}(x_k) + \al)) \chi_{[y_k
\neq \tilde f(x_k)]} \\
&&+ \sum_{k=1}^N g(y_kf^{(j)}(x_k) + \al)
\end{eqnarray*}
This shows that, for a given $\al$, the classifier $\tilde f$ must
minimize a weighted empirical error with weights
$$ w_k  = g(y_kf^{(j)}(x_k) - \al) - g(y_kf^{(j)}(x_k) + \al) \geq 0,$$
since $g$ is decreasing. 
However, the combined derivation of $\al_{j+1}$ and $f_{j+1}$ is, in general, computationally intensive, typically requiring to  loop over an estimation of
$f_{j+1}$ with fixed $\al_{j+1}$, and optimization of $\al_{j+1}$
by gradient descent, for example, with fixed $f_{j+1}$. There is anyway one situation in which the
computation can be made analytically: it is when the function $g$ is
exponential, $g(t) = e^{- t}$. In this case, the weights are
$$w_k = (e^{\al} - e^{-\al})   e^{- y_k f^{(j)}(x_k)}$$
and the constant factor $(e^{\al} - e^{-\al})$ does not affect
the solution of the minimization in $\tilde f$, which can therefore be
computed with
$w_k = \exp(- y_k f^{(j)}(x_k))$, a problem which is independent of $\al$. The value of $\al$, afterward,
can be computed by  minimizing:
$$
(e^{\al} - e^{-\al}) \left(\sum_k w_k \chi_{[y_k
\neq \tilde f(x_k)]}\right) + e^{-\al} \sum_k w_k
$$
This yields $ \al = \frac{1}{2} \log((1-E)/E)$ with
$$
E = \sum_k w_k \chi_{[y_k
\neq \tilde f(x_k)]}/ \sum_k w_k.
$$
This provides the previously described algorithm.

\bigskip
One critic that can be addressed to this algorithm is that it does
not properly handles the maximization of the margin, since the
$f^{(j)}$'s  are not normalized. One possibility is to require that
the $\al_q$ are positive and sum to 1.
The new classifier $f^{(j+1)}$ will then have the form $(1-\al_{j+1})f^{(j)} + \al_{j+1}
f_{j+1}$, and the function to minimize would be
$$
U(\al, \tilde f) = \sum_{k=1}^N g((1-\al) y_kf^{(j)}(x_k) + \al y_k \tilde
f(x_k)).
$$

Taking as before an exponential $g$: $g(t) = e^{-\la t}$, the final
problem becomes less simple because $\al$ and $\tilde f$ do not
separate: the weights are    
$$w_k = (e^{\la \al} - e^{-\la \al})   e^{-\la(1-\al) y_k f^{(j)}(x_k)}$$
and $\al$ cannot be factorized out. The solution has to be computed
numerically, which can be costly in some cases. The general algorithm
for creating a new component $f_{j+1}$ is to loop over the following
two steps, starting with an initial guess $f_{j+1}$ (computed, for
example, with the weights provided by the standard boosting algorithm).
\begin{enumerate}[label=(\arabic*),wide]
\item Minimize $U$ with respect to $\al$.
\item Update $f_{j+1}$ using the new weights 
$$
w_k = g((1-\al) y_kf^{(j)}(x_k) - \al) - g((1-\al) y_kf^{(j)}(x_k) + \al)
$$
\end{enumerate}
One can compute the derivatives of $U$ with respect to $\al$; we have
\begin{eqnarray*}
\prt_\al{U} & = & \sum_{k=1}^N y_k (f(x_k)-f^{(j)}(x_k)) g'((1-\al) y_kf^{(j)}(x_k) + \al y_k \tilde
f(x_k))\\
{\prt^2_\al U} & = & \sum_{k=1}^N (f(x_k)-f^{(j)}(x_k))^2 g''((1-\al) y_kf^{(j)}(x_k) + \al y_k \tilde
f(x_k))
\end{eqnarray*}
One sees in particular that $U$ is convex as a function of $\al$ as
soon as $g$ is convex (as is the exponential). A typical Newton-Raphson algorithm to solve
step 1 is
$$
\alpha(t+1) = \alpha(t) - \frac{1}{2} \prt_\al {U}/
{\prt^2_\al U}.
$$


\subsection{Choice of $f_j$}
The ``simple classifiers'' $f_j$ can be anything that accepts
weights. Notice that, for the algorithm to work, the
classifier does not have to minimize the weighted  classification error, but
simply to reduce it: the cost in the Adaboost algorithm is 
\begin{eqnarray*}
\CE{(j+1)}  &=& (e^{\al_{j+1}} - e^{-\al_{j+1}}) \sum_{k=1}^N w_k^{(j+1)} \chi_{[y_k
\neq \tilde f_{j+1}(x_k)]} + e^{-\al_{j+1}} \CE^{(j)}\\
&=&  \left((e^{\al_{j+1}} - e^{-\al_{j+1}}) \ep_{j+1} +
e^{-\al_{j+1}}\right)\CE^{(j)}
\end{eqnarray*}
where $\ep_{j+1}$ is the error corresponding to the $(j+1)$th
classifier alone and $\CE^{(j+1)}$ the combined cost. So, to have a
reduction of the cost, it suffices that 
$$(e^{\al_{j+1}} - e^{-\al_{j+1}}) \ep_{j+1} +
e^{-\al_{j+1}} < 1\,.$$
This is exactly the expression minimized by the boosting algorithm, and
at the optimal, $\al_{j+1} = \frac{1}{2}
\log((1-\ep_{j+1})/\ep_{j+1})$, this factor is $2
\sqrt{\ep_{j+1}(1-\ep_{j+1})}$, which is always smaller than 1, and
strictly smaller as soon as the error is strictly smaller than
$1/2$. So, exact minimization of the error is not needed at each step.

This implies that estimators such as LDA, linear separating hyperplanes,
etc. can be used as soon as they do better than chance. They do not
directly minimize the error, but weights on the training set always
have a natural way to be incorporated in their learning algorithms. For
example, assuming that weights sum to 1, a weighted LDA would use the class
centroids
$$
\mu_g = \sum_{k, y_k=g} w_k x_k
$$
and within-class covariance matrix
$$
\sum_{k=1}^N w_k (x_k - \mu_{y_k})(x_k - \mu_{y_k})^T.
$$ 

Similarly, a weighted SVM would minimize
$$
\norm{\be}^2 + \lambda N \sum_{k=1}^N w_k \xi_k
$$
under the constraints $\xi_k \geq 0$, $y_k(\be_0 + \scp{\be}{x_k})
\geq 1 - \xi_k$. In the dual variables, the problem becomes: maximize
$$
\sum_{k=1}^N \al_k - \frac{1}{2} \sum_{k,l=1}^N \al_k\al_l y_ky_l
K(x_k, x_l)
$$
under the constraints $0\leq \al_k\leq N\ga w_k$, $\sum_{k=1}^N \al_k y_k =
0$. One retrieves the usual SVM formulation with the change of
variables $\al_k \to \al_k/w_k$ and $y_k \to w_ky_k$.

\subsection{Boosting and regression}
The same idea can be applied, up to some extent, to non-linear regression problems. Here,
the coefficients  $\al_j$ are not needed because, most of the time,
if a function $f$ belongs to a regression class, the function $\al f$
will also belong to it. So, in this case, the issue is to find 
$$
f = f_1 + \cdots + f_K$$
by a greedy (online) procedure, to minimize, say, the mean square
error. Then, the procedure at step $j+1$ is to minimize
$$
\sum_{k=1}^N (y_k - f^{(j)}(x_k) - f(x_k))^2,
$$
which corresponds to the regression of the residuals $y_k -
f^{(j)}(x_k)$ on $x_k$. This has been applied with tree models, which are described in the next chapter.

